% GENERATED FILE. Edit canonical lessons, then run npm run generate:books.
% canonical-source-hash: fnv1a64:8591d56dd625fee1
% canonical-lessons: FR-C75-sol, FR-C75-escalier, FR-C75-miroir, FR-C75-salon, FR-C75-canape

\chapter{Floor, Stairs, Mirror, Living room, Sofa}
\label{ch:fr-floor-stairs-mirror-living-room-sofa}
\hlchaptermodality{75}

\begin{chapteropening}
\textbf{By the end of this chapter:} \emph{I can name the floor, the stairs, a mirror, a living room, a sofa.}

Complete the last lesson of chapter 75: I can name the floor, the stairs, a mirror, a living room, a sofa.
\end{chapteropening}

\section[le sol]{le sol — the floor, the ground}

\label{lesson:FR-C75-sol}



\begin{quote}
Before the new one: say the French for a bathroom, then the French for a bedroom.
\end{quote}

\subsection*{You'll want to know: le sol}
\textbf{le sol} — \textquotedblleft{}the floor, the ground\textquotedblright{}.

\begin{cousinweb}
\textbf{sol} $\leftarrow$ Latin \emph{solum}, \textquotedblleft{}ground, floor\textquotedblright{} — a learned borrowing. (EtymDB 2.1)
\end{cousinweb}

\begin{grammarlens}[title={What you've built}]
One more word for this chapter.
\end{grammarlens}

\subsection*{Your turn}
\begin{itemize}
\raggedright
  \item \emph{Say it:} \emph{le sol}
  \item \emph{Say it:} \emph{le sol}, once more
  \item \emph{Say it:} say \emph{le sol}
  \item \emph{Recall:} say the French for a bathroom, then the French for a bedroom, then say \emph{le sol} again
\end{itemize}

\begin{tcolorbox}[breakable,colback=teal!4,colframe=teal!35!black,title={Before you move on}]
What does le sol mean? (\textquotedblleft{}The floor, the ground\textquotedblright{}.) Say it once more.
\end{tcolorbox}
\section[l'escalier]{l'escalier — the stairs}

\label{lesson:FR-C75-escalier}



\begin{quote}
Before the new one: say the French for a bedroom, then the French for the floor.
\end{quote}

\subsection*{You'll want to know: l'escalier}
\textbf{l'escalier} — \textquotedblleft{}the stairs\textquotedblright{}.

\begin{cousinweb}
\textbf{escalier} $\leftarrow$ Occitan \emph{escalier} $\leftarrow$ Latin \emph{scalaria}, \textquotedblleft{}steps\textquotedblright{}, from \emph{scala}, \textquotedblleft{}a ladder\textquotedblright{} — borrowed from the south in the 1500s. (EtymDB 2.1 has the Middle French form.) The word French inherited from \emph{scala} is the one for \textquotedblleft{}ladder\textquotedblright{}.
\end{cousinweb}

\begin{grammarlens}[title={What you've built}]
One more word for this chapter.
\end{grammarlens}

\subsection*{Your turn}
\begin{itemize}
\raggedright
  \item \emph{Say it:} \emph{l'escalier}
  \item \emph{Say it:} \emph{l'escalier}, once more
  \item \emph{Say it:} say \emph{l'escalier}
  \item \emph{Recall:} say the French for a bedroom, then the French for the floor, then say \emph{l'escalier} again
\end{itemize}

\begin{tcolorbox}[breakable,colback=teal!4,colframe=teal!35!black,title={Before you move on}]
What does l'escalier mean? (\textquotedblleft{}The stairs\textquotedblright{}.) Say it once more.
\end{tcolorbox}
\section[le miroir]{le miroir — a mirror}

\label{lesson:FR-C75-miroir}



\begin{quote}
Before the new one: say the French for the floor, then the French for the stairs.
\end{quote}

\subsection*{You'll want to know: le miroir}
\textbf{le miroir} — \textquotedblleft{}a mirror\textquotedblright{}.

\begin{cousinweb}
\textbf{miroir} — built in Old French from the verb \emph{mirer}, \textquotedblleft{}to look at\textquotedblright{}, $\leftarrow$ Latin \emph{mirari}, \textquotedblleft{}to wonder at\textquotedblright{}. (EtymDB 2.1)
\end{cousinweb}

\begin{grammarlens}[title={What you've built}]
One more word for this chapter.
\end{grammarlens}

\subsection*{Your turn}
\begin{itemize}
\raggedright
  \item \emph{Say it:} \emph{le miroir}
  \item \emph{Say it:} \emph{le miroir}, once more
  \item \emph{Say it:} say \emph{le miroir}
  \item \emph{Recall:} say the French for the floor, then the French for the stairs, then say \emph{le miroir} again
\end{itemize}

\begin{tcolorbox}[breakable,colback=teal!4,colframe=teal!35!black,title={Before you move on}]
What does le miroir mean? (\textquotedblleft{}A mirror\textquotedblright{}.) Say it once more.
\end{tcolorbox}
\section[le salon]{le salon — a living room}

\label{lesson:FR-C75-salon}



\begin{quote}
Before the new one: say the French for the stairs, then the French for a mirror.
\end{quote}

\subsection*{You'll want to know: le salon}
\textbf{le salon} — \textquotedblleft{}a living room\textquotedblright{}.

\begin{cousinweb}
\textbf{salon} $\leftarrow$ Italian \emph{salone}, \textquotedblleft{}a large hall\textquotedblright{}, from \emph{sala}, \textquotedblleft{}a hall\textquotedblright{}, a Germanic word — borrowed. (EtymDB 2.1)
\end{cousinweb}

\begin{grammarlens}[title={What you've built}]
One more word for this chapter.
\end{grammarlens}

\subsection*{Your turn}
\begin{itemize}
\raggedright
  \item \emph{Say it:} \emph{le salon}
  \item \emph{Say it:} \emph{le salon}, once more
  \item \emph{Say it:} say \emph{le salon}
  \item \emph{Recall:} say the French for the stairs, then the French for a mirror, then say \emph{le salon} again
\end{itemize}

\begin{tcolorbox}[breakable,colback=teal!4,colframe=teal!35!black,title={Before you move on}]
What does le salon mean? (\textquotedblleft{}A living room\textquotedblright{}.) Say it once more.
\end{tcolorbox}
\section[le canapé]{le canapé — a sofa}

\label{lesson:FR-C75-canape}



\begin{quote}
Before the new one: say the French for a mirror, then the French for a living room.
\end{quote}

\subsection*{You'll want to know: le canapé}
\textbf{le canapé} — \textquotedblleft{}a sofa\textquotedblright{}.

\begin{cousinweb}
\textbf{canapé} $\leftarrow$ Medieval Latin \emph{canapeum} $\leftarrow$ Latin \emph{conopeum}, \textquotedblleft{}a bed with a mosquito net\textquotedblright{}, from Greek \emph{konopeion}, built on \emph{konops}, \textquotedblleft{}mosquito\textquotedblright{} — a learned borrowing. (EtymDB 2.1)
\end{cousinweb}

\begin{grammarlens}[title={What you've built}]
One more word for this chapter.
\end{grammarlens}

\subsection*{Your turn}
\begin{itemize}
\raggedright
  \item \emph{Say it:} \emph{le canapé}
  \item \emph{Say it:} \emph{le canapé}, once more
  \item \emph{Say it:} say \emph{le canapé}
  \item \emph{Recall:} say the French for a mirror, then the French for a living room, then say \emph{le canapé} again
\end{itemize}

\begin{tcolorbox}[breakable,colback=teal!4,colframe=teal!35!black,title={Before you move on}]
What does le canapé mean? (\textquotedblleft{}A sofa\textquotedblright{}.) Say it once more.
\end{tcolorbox}
